\chapter{Plan of Battle}

\lettrine{T}{he} night was already far advanced when the Abbé Fouquet joined his brother. Gourville had accompanied him. These three men, pale with dread of future events, resembled less three powers of the day than three conspirators, united by one single thought of violence. Fouquet walked for a long time, with his eyes fixed upon the floor, striking his hands one against the other. At length, taking courage, in the midst of a deep sigh: <Abbé,> said he, <you were speaking to me only to-day of certain people you maintain.>

<Yes, monsieur,> replied the Abbé.

<Tell me precisely who are these people.> The Abbé hesitated.

<Come! no fear, I am not threatening; no romancing, for I am not joking.>

<Since you demand the truth, monseigneur, here it is:—I have a hundred and twenty friends or companions of pleasure, who are sworn to me as the thief is to the gallows.>

<And you think you can depend on them?>

<Entirely.>

<And you will not compromise yourself?>

<I will not even make my appearance.>

<Are they men of resolution?>

<They would burn Paris, if I promised them they should not be burnt in turn.>

<The thing I ask of you, Abbé,> said Fouquet, wiping the sweat which fell from his brow, <is to throw your hundred and twenty men upon the people I will point out to you, at a certain moment given—is it possible?>

<It will not be the first time such a thing has happened to them, monseigneur.>

<That is well: but would these bandits attack an armed force?>

<They are used to that.>

<Then get your hundred and twenty men together, Abbé.>

<Directly. But where?>

<On the road to Vincennes, to-morrow, at two o'clock precisely.>

<To carry off Lyodot and D'Eymeris? There will be blows to be got!>

<A number, no doubt; are you afraid?>

<Not for myself, but for you.>

<Your men will know, then, what they have to do?>

<They are too intelligent not to guess it. Now, a minister who gets up a riot against his king—exposes himself\longdash>

<Of what importance is that to you, I pray? Besides, if I fall, you fall with me.>

<It would then be more prudent, monsieur, not to stir in the affair, and leave the king to take this little satisfaction.>

<Think well of this, Abbé, Lyodot and D'Eymeris at Vincennes are a prelude of ruin for my house. I repeat it—I arrested, you will be imprisoned—I imprisoned, you will be exiled.>

<Monsieur, I am at your orders; have you any to give me?>

<What I told you—I wish that, to-morrow, the two financiers of whom they mean to make victims, whilst there remain so many criminals unpunished, should be snatched from the fury of my enemies. Take your measures accordingly. Is it possible?>

<It is possible.>

<Describe your plan.>

<It is of rich simplicity. The ordinary guard at executions consists of twelve archers.>

<There will be a hundred to-morrow.>

<I reckon so. I even say more—there will be two hundred.>

<Then your hundred and twenty men will not be enough.>

<Pardon me. In every crowd composed of a hundred thousand spectators, there are ten thousand bandits or cut-purses—only they dare not take the initiative.>

<Well?>

<There will then be, to-morrow, on the Place de Grève, which I choose as my battle-field, ten thousand auxiliaries to my hundred and twenty men. The attack commenced by the latter, the others will finish it.>

<That all appears feasible. But what will be done with regard to the prisoners upon the Place de Grève?>

<This: they must be thrust into some house—that will make a siege necessary to get them out again. And stop! here is another idea, more sublime still: certain houses have two issues—one upon the Place, and the other into the Rue de la Mortellerie, or la Vannerie, or la Tixeranderie. The prisoners entering by one door will go out at another.>

<Yes; but fix upon something positive.>

<I am seeking to do so.>

<And I,> cried Fouquet, <I have found it. Listen to what has occurred to me at this moment.>

<I am listening.>

Fouquet made a sign to Gourville, who appeared to understand. <One of my friends lends me sometimes the keys of a house which he rents, Rue Baudoyer, the spacious gardens of which extend behind a certain house on the Place de Grève.>

<That is the place for us,> said the Abbé. <What house?>

<A cabaret, pretty well frequented, whose sign represents the image of Notre Dame.>

<I know it,> said the Abbé.

<This cabaret has windows opening upon the Place, a place of exit into the court, which must abut upon the gardens of my friend by a door of communication.>

<Good!> said the Abbé.

<Enter by the cabaret, take the prisoners in; defend the door while you enable them to fly by the garden and the Place Baudoyer.>

<That is all plain. Monsieur, you would make an excellent general, like monsieur le prince.>

<Have you understood me?>

<Perfectly well.>

<How much will it amount to, to make your bandits all drunk with wine, and to satisfy them with gold?>

<Oh, monsieur, what an expression! Oh! monsieur, if they heard you! some of them are very susceptible.>

<I mean to say they must be brought to the point where they cannot tell the heavens from the earth; for I shall to-morrow contend with the king; and when I fight I mean to conquer—please to understand.>

<It shall be done, monsieur. Give me your other ideas.>

<That is your business.>

<Then give me your purse.>

<Gourville, count a hundred thousand livres for the Abbé.>

<Good! and spare nothing, did you not say?>

<Nothing.>

<That is well.>

<Monseigneur,> objected Gourville, <if this should be known, we should lose our heads.>

<Eh! Gourville,> replied Fouquet, purple with anger, <you excite my pity. Speak for yourself, if you please. My head does not shake in that manner upon my shoulders. Now, Abbé, is everything arranged?>

<Everything.>

<At two o'clock to-morrow.>

<At twelve, because it will be necessary to prepare our auxiliaries in a secret manner.>

<That is true; do not spare the wine of the cabaretier.>

<I will spare neither his wine nor his house,> replied the Abbé, with a sneering laugh. <I have my plan, I tell you; leave me to set it in operation, and you shall see.>

<Where shall you be yourself?>

<Everywhere; nowhere.>

<And how shall I receive information?>

<By a courier whose horse shall be kept in the very same garden of your friend. À propos, the name of your friend?>

Fouquet looked again at Gourville. The latter came to the succour of his master, saying, [<The name is of no importance.>

Fouquet continued, <Accompany] monsieur l'Abbé, for several reasons, but the house is easily to be known—the \buildingname{Image-de-Notre-Dame} in the front, a garden, the only one in the quarter, behind.>\footnote{The text in the print copy is corrupt at this point. The suggested reading, in brackets, is by John Bursey.}

<Good, good! I will go and give notice to my soldiers.>

<Accompany him, Gourville,> said Fouquet, <and count him down the money. One moment, Abbé—one moment, Gourville—what name will be given to this carrying off?>

<A very natural one, monsieur—the Riot.>

<The riot on account of what? For, if ever the people of Paris are disposed to pay their court to the king, it is when he hangs financiers.>

<I will manage that,> said the Abbé.

<Yes; but you may manage it badly, and people will guess.>

<Not at all,—not at all. I have another idea.>

<What is that?>

<My men shall cry out, <Colbert, vive Colbert!> and shall throw themselves upon the prisoners as if they would tear them in pieces, and shall force them from the gibbets, as too mild a punishment.>

<Ah! that is an idea,> said Gourville. <\swear{Peste!} monsieur l'Abbé, what an imagination you have!>

<Monsieur, we are worthy of our family,> replied the Abbé, proudly.

<Strange fellow,> murmured Fouquet. Then he added, <That is ingenious. Carry it out, but shed no blood.>

Gourville and the Abbé set off together, with their heads full of the meditated riot. The superintendent laid himself down upon some cushions, half valiant with respect to the sinister projects of the morrow, half dreaming of love. 